\documentclass[11pt,letterpaper]{article}
\usepackage[technical-reference]{cornell-notes}
\usepackage{needspace}

\title{Clause 7: Expressions - Cornell Notes}
\author{}
\date{}
\setDocTitle{Clause 7: Expressions}
\setDocSubtitle{C++ 2024 Cornell Notes}
\setDocOwner{C++ 2024 Cornell Notes}
\setDocAuthor{}
\setDocDate{}
\setCornellCollection{C++ 2024 Cornell Notes}
\setCornellUnitType{Clause}
\setCornellUnitNumber{7}
\setCornellUnitTitle{Expressions}
\setCornellCompanionLabel{C++ Technical Reference}
\hypersetup{pdftitle={Clause 7: Expressions - Cornell Notes},pdfsubject={C++ technical study notes},pdfauthor={}}

\begin{document}
\maketitle
\begin{CornellReferenceOrientationBox}
\textbf{Purpose.} Defines value categories, types, conversions, operators, evaluation semantics, temporary materialization, allocation expressions, casts, lambdas, requires-expressions, and constant evaluation.
\par\textbf{Role.} Specifies how computations are formed, typed, sequenced, and evaluated.
\par\textbf{Principal dependencies.} Uses names, types, objects, lifetimes, and execution rules; feeds statements, initialization, overload resolution, and templates.
\par\textbf{Primary audience.} programmers, compiler front-end implementers, and correctness tools.
\par\textbf{Major subject groups.} Preamble; Properties of expressions; Standard conversions; Usual arithmetic conversions; Primary expressions; Compound expressions; Constant expressions.
\end{CornellReferenceOrientationBox}
\section{Learning Objectives}
\begin{itemize}[label=\textcolor{CornellReferenceTeal}{$\blacktriangleright$}]
\item Define the governing ideas of Preamble and state the consequences for a program or implementation.
\item Identify the governing ideas of Properties of expressions and state the consequences for a program or implementation.
\item Distinguish the governing ideas of Standard conversions and state the consequences for a program or implementation.
\item Explain the governing ideas of Usual arithmetic conversions and state the consequences for a program or implementation.
\item Trace the governing ideas of Primary expressions and state the consequences for a program or implementation.
\item Apply the governing ideas of Compound expressions and state the consequences for a program or implementation.
\item Analyze the governing ideas of Constant expressions and state the consequences for a program or implementation.
\item Evaluate the governing ideas of Value category and state the consequences for a program or implementation.
\item Diagnose the governing ideas of Type and state the consequences for a program or implementation.
\item Compare the governing ideas of Context dependence and state the consequences for a program or implementation.
\item Verify the governing ideas of General and state the consequences for a program or implementation.
\item Relate the governing ideas of Lvalue-to-rvalue conversion and state the consequences for a program or implementation.
\item Classify the governing ideas of Array-to-pointer conversion and state the consequences for a program or implementation.
\item Justify the governing ideas of Function-to-pointer conversion and state the consequences for a program or implementation.
\item Audit the governing ideas of Temporary materialization conversion and state the consequences for a program or implementation.
\item Summarize the governing ideas of Qualification conversions and state the consequences for a program or implementation.
\end{itemize}
\section{Normative Structure Map}
\small\begin{longtable}{@{}>{\RaggedRight\arraybackslash}p{0.13\textwidth}>{\RaggedRight\arraybackslash}p{0.27\textwidth}>{\RaggedRight\arraybackslash}p{0.19\textwidth}>{\RaggedRight\arraybackslash}p{0.31\textwidth}@{}}
\toprule\textbf{Subclause} & \textbf{Subject} & \textbf{Rule category} & \textbf{Key dependencies / entities}\\\midrule\endfirsthead
\toprule\textbf{Subclause} & \textbf{Subject} & \textbf{Rule category} & \textbf{Key dependencies / entities}\\\midrule\endhead
\bottomrule\endfoot
7.1 & Preamble & core-language semantics & Uses names, types, objects, lifetimes, and execution rules; feeds statements, initialization, overload resolution, and templates.\\
7.2 & Properties of expressions & core-language semantics & Value category, Type, Context dependence\\
7.3 & Standard conversions & core-language semantics & General, Lvalue-to-rvalue conversion, Array-to-pointer conversion, Function-to-pointer conversion, and related subclauses\\
7.4 & Usual arithmetic conversions & core-language semantics & Uses names, types, objects, lifetimes, and execution rules; feeds statements, initialization, overload resolution, and templates.\\
7.5 & Primary expressions & core-language semantics & Grammar, Literals, This, Parentheses, and related subclauses\\
7.6 & Compound expressions & core-language semantics & Postfix expressions, Unary expressions, Explicit type conversion (cast notation), Pointer-to-member operators, and related subclauses\\
7.7 & Constant expressions & core-language semantics & Uses names, types, objects, lifetimes, and execution rules; feeds statements, initialization, overload resolution, and templates.\\
\end{longtable}\normalsize
\begin{CornellReferenceNormativeBox}A heading maps the clause; it does not replace the detailed conditions. For each operation, read requirements, effects, result, exceptions, complexity, remarks, and notes with their stated normative force.\end{CornellReferenceNormativeBox}
\subsection*{Rule Classification Guide}
\small\begin{longtable}{@{}>{\RaggedRight\arraybackslash}p{0.25\textwidth}>{\RaggedRight\arraybackslash}p{0.67\textwidth}@{}}\toprule\textbf{Label or category} & \textbf{How to read it}\\\midrule\endfirsthead\toprule\textbf{Label or category} & \textbf{How to read it (continued)}\\\midrule\endhead\bottomrule\endfoot
Well-formed / ill-formed & Formation is governed by syntax and semantic rules; an ill-formed program does not become well-formed because one implementation accepts it.\\
Diagnosable rule & A violation requires at least one diagnostic unless the wording places the case outside the diagnosable set.\\
No diagnostic required & The program can be ill-formed while the implementation has no diagnostic obligation.\\
Undefined behavior & No execution requirements are imposed; translation success does not establish safety or portability.\\
Unspecified behavior & One of the permitted outcomes may be chosen without documentation.\\
Implementation-defined behavior & The implementation chooses and documents the behavior.\\
Conditionally supported & The implementation may omit support but documents unsupported constructs.\\
\end{longtable}\normalsize
\section{Cornell Cue-and-Notes}
\Needspace{8\baselineskip}
\subsection*{7.1 Preamble}
\begin{longtable}{@{}>{\RaggedRight\arraybackslash\columncolor{CornellReferenceCueGray}}p{0.28\textwidth}>{\RaggedRight\arraybackslash}p{0.64\textwidth}@{}}
\rowcolor{CornellReferenceNavy}\color{white}\textbf{Cue / recall prompt} & \color{white}\textbf{Notes}\\\endfirsthead
\rowcolor{CornellReferenceNavy}\color{white}\textbf{Cue / recall prompt} & \color{white}\textbf{Notes (continued)}\\\endhead
\arrayrulecolor{CornellReferenceLineGray}
\textbf{What rule applies to Preamble?}\\[-1mm]{\scriptsize\CornellClauseRef{7.1}} & Frames the terminology, applicability, and relationships needed to interpret Preamble; detailed rules remain in the following subclauses.\\\hline
\end{longtable}
\begin{CornellReferenceSummaryBox}Preamble supplies the core-language semantics portion of this clause. The key review move is to connect its local wording to the clause-wide model, then classify each condition by normative force before relying on it.\end{CornellReferenceSummaryBox}
\Needspace{8\baselineskip}
\subsection*{7.2 Properties of expressions}
\begin{longtable}{@{}>{\RaggedRight\arraybackslash\columncolor{CornellReferenceCueGray}}p{0.28\textwidth}>{\RaggedRight\arraybackslash}p{0.64\textwidth}@{}}
\rowcolor{CornellReferenceNavy}\color{white}\textbf{Cue / recall prompt} & \color{white}\textbf{Notes}\\\endfirsthead
\rowcolor{CornellReferenceNavy}\color{white}\textbf{Cue / recall prompt} & \color{white}\textbf{Notes (continued)}\\\endhead
\arrayrulecolor{CornellReferenceLineGray}
\textbf{What rule applies to Value category?}\\[-1mm]{\scriptsize\CornellClauseRef{7.2.1}} & Classifies expressions in Value category by identity and value production, which controls materialization, reference binding, and overload behavior.\\\hline
\textbf{What rule applies to Type?}\\[-1mm]{\scriptsize\CornellClauseRef{7.2.2}} & Defines the core-language rules for Type, including formation, semantic effect, interactions with type and lookup, and any ill-formed or implementation-dependent cases.\\\hline
\textbf{What rule applies to Context dependence?}\\[-1mm]{\scriptsize\CornellClauseRef{7.2.3}} & Defines the core-language rules for Context dependence, including formation, semantic effect, interactions with type and lookup, and any ill-formed or implementation-dependent cases.\\\hline
\end{longtable}
\begin{CornellReferenceSummaryBox}Properties of expressions supplies the core-language semantics portion of this clause. The key review move is to connect its local wording to the clause-wide model, then classify each condition by normative force before relying on it.\end{CornellReferenceSummaryBox}
\Needspace{8\baselineskip}
\subsection*{7.3 Standard conversions}
\begin{longtable}{@{}>{\RaggedRight\arraybackslash\columncolor{CornellReferenceCueGray}}p{0.28\textwidth}>{\RaggedRight\arraybackslash}p{0.64\textwidth}@{}}
\rowcolor{CornellReferenceNavy}\color{white}\textbf{Cue / recall prompt} & \color{white}\textbf{Notes}\\\endfirsthead
\rowcolor{CornellReferenceNavy}\color{white}\textbf{Cue / recall prompt} & \color{white}\textbf{Notes (continued)}\\\endhead
\arrayrulecolor{CornellReferenceLineGray}
\textbf{What rule applies to General?}\\[-1mm]{\scriptsize\CornellClauseRef{7.3.1}} & Frames the terminology, applicability, and relationships needed to interpret General; detailed rules remain in the following subclauses.\\\hline
\textbf{When is Lvalue-to-rvalue conversion permitted?}\\[-1mm]{\scriptsize\CornellClauseRef{7.3.2}} & Specifies source and destination conditions for Lvalue-to-rvalue conversion, the resulting type or value category, and any information-loss, pointer, qualification, or lifetime consequence.\\\hline
\textbf{When is Array-to-pointer conversion permitted?}\\[-1mm]{\scriptsize\CornellClauseRef{7.3.3}} & Specifies source and destination conditions for Array-to-pointer conversion, the resulting type or value category, and any information-loss, pointer, qualification, or lifetime consequence.\\\hline
\textbf{When is Function-to-pointer conversion permitted?}\\[-1mm]{\scriptsize\CornellClauseRef{7.3.4}} & Specifies source and destination conditions for Function-to-pointer conversion, the resulting type or value category, and any information-loss, pointer, qualification, or lifetime consequence.\\\hline
\textbf{When is Temporary materialization conversion permitted?}\\[-1mm]{\scriptsize\CornellClauseRef{7.3.5}} & Specifies source and destination conditions for Temporary materialization conversion, the resulting type or value category, and any information-loss, pointer, qualification, or lifetime consequence.\\\hline
\textbf{When is Qualification conversions permitted?}\\[-1mm]{\scriptsize\CornellClauseRef{7.3.6}} & Specifies source and destination conditions for Qualification conversions, the resulting type or value category, and any information-loss, pointer, qualification, or lifetime consequence.\\\hline
\textbf{What rule applies to Integral promotions?}\\[-1mm]{\scriptsize\CornellClauseRef{7.3.7}} & Specifies source and destination conditions for Integral promotions, the resulting type or value category, and any information-loss, pointer, qualification, or lifetime consequence.\\\hline
\textbf{What rule applies to Floating-point promotion?}\\[-1mm]{\scriptsize\CornellClauseRef{7.3.8}} & Specifies source and destination conditions for Floating-point promotion, the resulting type or value category, and any information-loss, pointer, qualification, or lifetime consequence.\\\hline
\textbf{When is Integral conversions permitted?}\\[-1mm]{\scriptsize\CornellClauseRef{7.3.9}} & Specifies source and destination conditions for Integral conversions, the resulting type or value category, and any information-loss, pointer, qualification, or lifetime consequence.\\\hline
\textbf{When is Floating-point conversions permitted?}\\[-1mm]{\scriptsize\CornellClauseRef{7.3.10}} & Specifies source and destination conditions for Floating-point conversions, the resulting type or value category, and any information-loss, pointer, qualification, or lifetime consequence.\\\hline
\textbf{When is Floating-integral conversions permitted?}\\[-1mm]{\scriptsize\CornellClauseRef{7.3.11}} & Specifies source and destination conditions for Floating-integral conversions, the resulting type or value category, and any information-loss, pointer, qualification, or lifetime consequence.\\\hline
\textbf{When is Pointer conversions permitted?}\\[-1mm]{\scriptsize\CornellClauseRef{7.3.12}} & Specifies source and destination conditions for Pointer conversions, the resulting type or value category, and any information-loss, pointer, qualification, or lifetime consequence.\\\hline
\textbf{When is Pointer-to-member conversions permitted?}\\[-1mm]{\scriptsize\CornellClauseRef{7.3.13}} & Specifies source and destination conditions for Pointer-to-member conversions, the resulting type or value category, and any information-loss, pointer, qualification, or lifetime consequence.\\\hline
\textbf{When is Function pointer conversions permitted?}\\[-1mm]{\scriptsize\CornellClauseRef{7.3.14}} & Specifies source and destination conditions for Function pointer conversions, the resulting type or value category, and any information-loss, pointer, qualification, or lifetime consequence.\\\hline
\textbf{When is Boolean conversions permitted?}\\[-1mm]{\scriptsize\CornellClauseRef{7.3.15}} & Specifies source and destination conditions for Boolean conversions, the resulting type or value category, and any information-loss, pointer, qualification, or lifetime consequence.\\\hline
\end{longtable}
\begin{CornellReferenceSummaryBox}Standard conversions supplies the core-language semantics portion of this clause. The key review move is to connect its local wording to the clause-wide model, then classify each condition by normative force before relying on it.\end{CornellReferenceSummaryBox}
\Needspace{8\baselineskip}
\subsection*{7.4 Usual arithmetic conversions}
\begin{longtable}{@{}>{\RaggedRight\arraybackslash\columncolor{CornellReferenceCueGray}}p{0.28\textwidth}>{\RaggedRight\arraybackslash}p{0.64\textwidth}@{}}
\rowcolor{CornellReferenceNavy}\color{white}\textbf{Cue / recall prompt} & \color{white}\textbf{Notes}\\\endfirsthead
\rowcolor{CornellReferenceNavy}\color{white}\textbf{Cue / recall prompt} & \color{white}\textbf{Notes (continued)}\\\endhead
\arrayrulecolor{CornellReferenceLineGray}
\textbf{When is Usual arithmetic conversions permitted?}\\[-1mm]{\scriptsize\CornellClauseRef{7.4}} & Specifies source and destination conditions for Usual arithmetic conversions, the resulting type or value category, and any information-loss, pointer, qualification, or lifetime consequence.\\\hline
\end{longtable}
\begin{CornellReferenceSummaryBox}Usual arithmetic conversions supplies the core-language semantics portion of this clause. The key review move is to connect its local wording to the clause-wide model, then classify each condition by normative force before relying on it.\end{CornellReferenceSummaryBox}
\Needspace{8\baselineskip}
\subsection*{7.5 Primary expressions}
\begin{longtable}{@{}>{\RaggedRight\arraybackslash\columncolor{CornellReferenceCueGray}}p{0.28\textwidth}>{\RaggedRight\arraybackslash}p{0.64\textwidth}@{}}
\rowcolor{CornellReferenceNavy}\color{white}\textbf{Cue / recall prompt} & \color{white}\textbf{Notes}\\\endfirsthead
\rowcolor{CornellReferenceNavy}\color{white}\textbf{Cue / recall prompt} & \color{white}\textbf{Notes (continued)}\\\endhead
\arrayrulecolor{CornellReferenceLineGray}
\textbf{What rule applies to Grammar?}\\[-1mm]{\scriptsize\CornellClauseRef{7.5.1}} & Defines the recognized forms for Grammar; syntactic acceptance does not replace the semantic constraints stated with the production.\\\hline
\textbf{What rule applies to Literals?}\\[-1mm]{\scriptsize\CornellClauseRef{7.5.2}} & Defines the core-language rules for Literals, including formation, semantic effect, interactions with type and lookup, and any ill-formed or implementation-dependent cases.\\\hline
\textbf{What rule applies to This?}\\[-1mm]{\scriptsize\CornellClauseRef{7.5.3}} & Defines the core-language rules for This, including formation, semantic effect, interactions with type and lookup, and any ill-formed or implementation-dependent cases.\\\hline
\textbf{What rule applies to Parentheses?}\\[-1mm]{\scriptsize\CornellClauseRef{7.5.4}} & Defines the core-language rules for Parentheses, including formation, semantic effect, interactions with type and lookup, and any ill-formed or implementation-dependent cases.\\\hline
\textbf{What rule applies to Names?}\\[-1mm]{\scriptsize\CornellClauseRef{7.5.5}} & Defines the core-language rules for Names, including formation, semantic effect, interactions with type and lookup, and any ill-formed or implementation-dependent cases.\\\hline
\textbf{What rule applies to Lambda expressions?}\\[-1mm]{\scriptsize\CornellClauseRef{7.5.6}} & Defines the core-language rules for Lambda expressions, including formation, semantic effect, interactions with type and lookup, and any ill-formed or implementation-dependent cases.\\\hline
\textbf{What rule applies to Fold expressions?}\\[-1mm]{\scriptsize\CornellClauseRef{7.5.7}} & Defines the core-language rules for Fold expressions, including formation, semantic effect, interactions with type and lookup, and any ill-formed or implementation-dependent cases.\\\hline
\textbf{What rule applies to Requires expressions?}\\[-1mm]{\scriptsize\CornellClauseRef{7.5.8}} & Defines the core-language rules for Requires expressions, including formation, semantic effect, interactions with type and lookup, and any ill-formed or implementation-dependent cases.\\\hline
\end{longtable}
\begin{CornellReferenceSummaryBox}Primary expressions supplies the core-language semantics portion of this clause. The key review move is to connect its local wording to the clause-wide model, then classify each condition by normative force before relying on it.\end{CornellReferenceSummaryBox}
\Needspace{5\baselineskip}\paragraph{Detailed coverage ledger.}
\begin{description}[style=nextline,leftmargin=2.8cm,font=\normalfont\bfseries\color{CornellReferenceTeal}]
\item[7.5.5.1] \textbf{General.} Frames the terminology, applicability, and relationships needed to interpret General; detailed rules remain in the following subclauses.
\item[7.5.5.2] \textbf{Unqualified names.} Defines the core-language rules for Unqualified names, including formation, semantic effect, interactions with type and lookup, and any ill-formed or implementation-dependent cases.
\item[7.5.5.3] \textbf{Qualified names.} Defines the core-language rules for Qualified names, including formation, semantic effect, interactions with type and lookup, and any ill-formed or implementation-dependent cases.
\item[7.5.5.4] \textbf{Destruction.} Defines the core-language rules for Destruction, including formation, semantic effect, interactions with type and lookup, and any ill-formed or implementation-dependent cases.
\item[7.5.6.1] \textbf{General.} Frames the terminology, applicability, and relationships needed to interpret General; detailed rules remain in the following subclauses.
\item[7.5.6.2] \textbf{Closure types.} Defines the core-language rules for Closure types, including formation, semantic effect, interactions with type and lookup, and any ill-formed or implementation-dependent cases.
\item[7.5.6.3] \textbf{Captures.} Defines the core-language rules for Captures, including formation, semantic effect, interactions with type and lookup, and any ill-formed or implementation-dependent cases.
\item[7.5.8.1] \textbf{General.} Frames the terminology, applicability, and relationships needed to interpret General; detailed rules remain in the following subclauses.
\item[7.5.8.2] \textbf{Simple requirements.} States the admissibility and semantic obligations for Simple requirements. Separate compile-time participation rules from caller preconditions and runtime effects.
\item[7.5.8.3] \textbf{Type requirements.} States the admissibility and semantic obligations for Type requirements. Separate compile-time participation rules from caller preconditions and runtime effects.
\item[7.5.8.4] \textbf{Compound requirements.} States the admissibility and semantic obligations for Compound requirements. Separate compile-time participation rules from caller preconditions and runtime effects.
\item[7.5.8.5] \textbf{Nested requirements.} States the admissibility and semantic obligations for Nested requirements. Separate compile-time participation rules from caller preconditions and runtime effects.
\end{description}
\Needspace{8\baselineskip}
\subsection*{7.6 Compound expressions}
\begin{longtable}{@{}>{\RaggedRight\arraybackslash\columncolor{CornellReferenceCueGray}}p{0.28\textwidth}>{\RaggedRight\arraybackslash}p{0.64\textwidth}@{}}
\rowcolor{CornellReferenceNavy}\color{white}\textbf{Cue / recall prompt} & \color{white}\textbf{Notes}\\\endfirsthead
\rowcolor{CornellReferenceNavy}\color{white}\textbf{Cue / recall prompt} & \color{white}\textbf{Notes (continued)}\\\endhead
\arrayrulecolor{CornellReferenceLineGray}
\textbf{What rule applies to Postfix expressions?}\\[-1mm]{\scriptsize\CornellClauseRef{7.6.1}} & Defines the core-language rules for Postfix expressions, including formation, semantic effect, interactions with type and lookup, and any ill-formed or implementation-dependent cases.\\\hline
\textbf{What rule applies to Unary expressions?}\\[-1mm]{\scriptsize\CornellClauseRef{7.6.2}} & Defines the core-language rules for Unary expressions, including formation, semantic effect, interactions with type and lookup, and any ill-formed or implementation-dependent cases.\\\hline
\textbf{When is Explicit type conversion (cast notation) permitted?}\\[-1mm]{\scriptsize\CornellClauseRef{7.6.3}} & Specifies source and destination conditions for Explicit type conversion (cast notation), the resulting type or value category, and any information-loss, pointer, qualification, or lifetime consequence.\\\hline
\textbf{What rule applies to Pointer-to-member operators?}\\[-1mm]{\scriptsize\CornellClauseRef{7.6.4}} & Defines the core-language rules for Pointer-to-member operators, including formation, semantic effect, interactions with type and lookup, and any ill-formed or implementation-dependent cases.\\\hline
\textbf{What rule applies to Multiplicative operators?}\\[-1mm]{\scriptsize\CornellClauseRef{7.6.5}} & Defines the core-language rules for Multiplicative operators, including formation, semantic effect, interactions with type and lookup, and any ill-formed or implementation-dependent cases.\\\hline
\textbf{What rule applies to Additive operators?}\\[-1mm]{\scriptsize\CornellClauseRef{7.6.6}} & Defines the core-language rules for Additive operators, including formation, semantic effect, interactions with type and lookup, and any ill-formed or implementation-dependent cases.\\\hline
\textbf{What rule applies to Shift operators?}\\[-1mm]{\scriptsize\CornellClauseRef{7.6.7}} & Defines the core-language rules for Shift operators, including formation, semantic effect, interactions with type and lookup, and any ill-formed or implementation-dependent cases.\\\hline
\textbf{What rule applies to Three-way comparison operator?}\\[-1mm]{\scriptsize\CornellClauseRef{7.6.8}} & Defines the core-language rules for Three-way comparison operator, including formation, semantic effect, interactions with type and lookup, and any ill-formed or implementation-dependent cases.\\\hline
\textbf{What rule applies to Relational operators?}\\[-1mm]{\scriptsize\CornellClauseRef{7.6.9}} & Defines the core-language rules for Relational operators, including formation, semantic effect, interactions with type and lookup, and any ill-formed or implementation-dependent cases.\\\hline
\textbf{What rule applies to Equality operators?}\\[-1mm]{\scriptsize\CornellClauseRef{7.6.10}} & Defines the core-language rules for Equality operators, including formation, semantic effect, interactions with type and lookup, and any ill-formed or implementation-dependent cases.\\\hline
\textbf{What rule applies to Bitwise AND operator?}\\[-1mm]{\scriptsize\CornellClauseRef{7.6.11}} & Defines the core-language rules for Bitwise AND operator, including formation, semantic effect, interactions with type and lookup, and any ill-formed or implementation-dependent cases.\\\hline
\textbf{What rule applies to Bitwise exclusive OR operator?}\\[-1mm]{\scriptsize\CornellClauseRef{7.6.12}} & Defines the core-language rules for Bitwise exclusive OR operator, including formation, semantic effect, interactions with type and lookup, and any ill-formed or implementation-dependent cases.\\\hline
\textbf{What rule applies to Bitwise inclusive OR operator?}\\[-1mm]{\scriptsize\CornellClauseRef{7.6.13}} & Defines the core-language rules for Bitwise inclusive OR operator, including formation, semantic effect, interactions with type and lookup, and any ill-formed or implementation-dependent cases.\\\hline
\textbf{What rule applies to Logical AND operator?}\\[-1mm]{\scriptsize\CornellClauseRef{7.6.14}} & Defines the core-language rules for Logical AND operator, including formation, semantic effect, interactions with type and lookup, and any ill-formed or implementation-dependent cases.\\\hline
\textbf{What rule applies to Logical OR operator?}\\[-1mm]{\scriptsize\CornellClauseRef{7.6.15}} & Defines the core-language rules for Logical OR operator, including formation, semantic effect, interactions with type and lookup, and any ill-formed or implementation-dependent cases.\\\hline
\textbf{What rule applies to Conditional operator?}\\[-1mm]{\scriptsize\CornellClauseRef{7.6.16}} & Defines the core-language rules for Conditional operator, including formation, semantic effect, interactions with type and lookup, and any ill-formed or implementation-dependent cases.\\\hline
\textbf{What rule applies to Yielding a value?}\\[-1mm]{\scriptsize\CornellClauseRef{7.6.17}} & Defines the core-language rules for Yielding a value, including formation, semantic effect, interactions with type and lookup, and any ill-formed or implementation-dependent cases.\\\hline
\textbf{What rule applies to Throwing an exception?}\\[-1mm]{\scriptsize\CornellClauseRef{7.6.18}} & Defines the failure model for Throwing an exception, including how an error is represented, reported, propagated, or converted into a diagnostic consequence.\\\hline
\textbf{What rule applies to Assignment and compound assignment operators?}\\[-1mm]{\scriptsize\CornellClauseRef{7.6.19}} & Governs state transfer or exchange for Assignment and compound assignment operators; check self-operation, validity after movement, exception behavior, and invalidation.\\\hline
\textbf{What rule applies to Comma operator?}\\[-1mm]{\scriptsize\CornellClauseRef{7.6.20}} & Defines the core-language rules for Comma operator, including formation, semantic effect, interactions with type and lookup, and any ill-formed or implementation-dependent cases.\\\hline
\end{longtable}
\begin{CornellReferenceSummaryBox}Compound expressions supplies the core-language semantics portion of this clause. The key review move is to connect its local wording to the clause-wide model, then classify each condition by normative force before relying on it.\end{CornellReferenceSummaryBox}
\Needspace{5\baselineskip}\paragraph{Detailed coverage ledger.}
\begin{description}[style=nextline,leftmargin=2.8cm,font=\normalfont\bfseries\color{CornellReferenceTeal}]
\item[7.6.1.1] \textbf{General.} Frames the terminology, applicability, and relationships needed to interpret General; detailed rules remain in the following subclauses.
\item[7.6.1.2] \textbf{Subscripting.} Defines the core-language rules for Subscripting, including formation, semantic effect, interactions with type and lookup, and any ill-formed or implementation-dependent cases.
\item[7.6.1.3] \textbf{Function call.} Defines the core-language rules for Function call, including formation, semantic effect, interactions with type and lookup, and any ill-formed or implementation-dependent cases.
\item[7.6.1.4] \textbf{Explicit type conversion (functional notation).} Specifies source and destination conditions for Explicit type conversion (functional notation), the resulting type or value category, and any information-loss, pointer, qualification, or lifetime consequence.
\item[7.6.1.5] \textbf{Class member access.} Defines the core-language rules for Class member access, including formation, semantic effect, interactions with type and lookup, and any ill-formed or implementation-dependent cases.
\item[7.6.1.6] \textbf{Increment and decrement.} Defines the core-language rules for Increment and decrement, including formation, semantic effect, interactions with type and lookup, and any ill-formed or implementation-dependent cases.
\item[7.6.1.7] \textbf{Dynamic cast.} Defines the core-language rules for Dynamic cast, including formation, semantic effect, interactions with type and lookup, and any ill-formed or implementation-dependent cases.
\item[7.6.1.8] \textbf{Type identification.} Defines the core-language rules for Type identification, including formation, semantic effect, interactions with type and lookup, and any ill-formed or implementation-dependent cases.
\item[7.6.1.9] \textbf{Static cast.} Defines the core-language rules for Static cast, including formation, semantic effect, interactions with type and lookup, and any ill-formed or implementation-dependent cases.
\item[7.6.1.10] \textbf{Reinterpret cast.} Defines the core-language rules for Reinterpret cast, including formation, semantic effect, interactions with type and lookup, and any ill-formed or implementation-dependent cases.
\item[7.6.1.11] \textbf{Const cast.} Defines the core-language rules for Const cast, including formation, semantic effect, interactions with type and lookup, and any ill-formed or implementation-dependent cases.
\item[7.6.2.1] \textbf{General.} Frames the terminology, applicability, and relationships needed to interpret General; detailed rules remain in the following subclauses.
\item[7.6.2.2] \textbf{Unary operators.} Defines the core-language rules for Unary operators, including formation, semantic effect, interactions with type and lookup, and any ill-formed or implementation-dependent cases.
\item[7.6.2.3] \textbf{Increment and decrement.} Defines the core-language rules for Increment and decrement, including formation, semantic effect, interactions with type and lookup, and any ill-formed or implementation-dependent cases.
\item[7.6.2.4] \textbf{Await.} Defines the core-language rules for Await, including formation, semantic effect, interactions with type and lookup, and any ill-formed or implementation-dependent cases.
\item[7.6.2.5] \textbf{Sizeof.} Defines the core-language rules for Sizeof, including formation, semantic effect, interactions with type and lookup, and any ill-formed or implementation-dependent cases.
\item[7.6.2.6] \textbf{Alignof.} Defines the core-language rules for Alignof, including formation, semantic effect, interactions with type and lookup, and any ill-formed or implementation-dependent cases.
\item[7.6.2.7] \textbf{noexcept operator.} Defines the core-language rules for noexcept operator, including formation, semantic effect, interactions with type and lookup, and any ill-formed or implementation-dependent cases.
\item[7.6.2.8] \textbf{New.} Defines the core-language rules for New, including formation, semantic effect, interactions with type and lookup, and any ill-formed or implementation-dependent cases.
\item[7.6.2.9] \textbf{Delete.} Defines the core-language rules for Delete, including formation, semantic effect, interactions with type and lookup, and any ill-formed or implementation-dependent cases.
\end{description}
\Needspace{8\baselineskip}
\subsection*{7.7 Constant expressions}
\begin{longtable}{@{}>{\RaggedRight\arraybackslash\columncolor{CornellReferenceCueGray}}p{0.28\textwidth}>{\RaggedRight\arraybackslash}p{0.64\textwidth}@{}}
\rowcolor{CornellReferenceNavy}\color{white}\textbf{Cue / recall prompt} & \color{white}\textbf{Notes}\\\endfirsthead
\rowcolor{CornellReferenceNavy}\color{white}\textbf{Cue / recall prompt} & \color{white}\textbf{Notes (continued)}\\\endhead
\arrayrulecolor{CornellReferenceLineGray}
\textbf{What rule applies to Constant expressions?}\\[-1mm]{\scriptsize\CornellClauseRef{7.7}} & Defines the core-language rules for Constant expressions, including formation, semantic effect, interactions with type and lookup, and any ill-formed or implementation-dependent cases.\\\hline
\end{longtable}
\begin{CornellReferenceSummaryBox}Constant expressions supplies the core-language semantics portion of this clause. The key review move is to connect its local wording to the clause-wide model, then classify each condition by normative force before relying on it.\end{CornellReferenceSummaryBox}
\section{Language-Lawyer and Library-Usage Pitfalls}
\begin{CornellReferencePitfallBox}\begin{itemize}[label=\textcolor{CornellReferenceAmber}{$\blacktriangleright$}]
\item Treating value category as a synonym for type.
\item Assuming operand evaluation order from operator spelling.
\item Believing a cast makes an invalid operation well-defined.
\item Ignoring temporary materialization and lifetime consequences.
\item Confusing contextual conversion with a general implicit conversion.
\item Treating a note, example, or recommended practice as though it imposed a mandatory program rule.
\end{itemize}\end{CornellReferencePitfallBox}
\section{Programmer and Implementation Checklist}
\begin{CornellReferenceChecklistBox}\begin{itemize}[label=$\square$]
\item Determine type and value category at every important subexpression.
\item List the conversions applied before the operator semantics.
\item Establish sequencing before relying on side effects.
\item Check pointer provenance, bounds, and lifetime.
\item Apply the precise cast restrictions rather than implementation behavior.
\item Verify constant-expression suitability in the actual context.
\item Include overload resolution and template dependence where relevant.
\item For \CornellClauseRef{7.1}, verify preamble against the precise rule category and all cited dependencies.
\item For \CornellClauseRef{7.2}, verify properties of expressions against the precise rule category and all cited dependencies.
\item For \CornellClauseRef{7.3}, verify standard conversions against the precise rule category and all cited dependencies.
\item For \CornellClauseRef{7.4}, verify usual arithmetic conversions against the precise rule category and all cited dependencies.
\end{itemize}\end{CornellReferenceChecklistBox}
\section{Clause Synthesis}
\begin{CornellReferenceOrientationBox}Defines value categories, types, conversions, operators, evaluation semantics, temporary materialization, allocation expressions, casts, lambdas, requires-expressions, and constant evaluation. Specifies how computations are formed, typed, sequenced, and evaluated. A sound review distinguishes syntax and participation from runtime preconditions, keeps notes and examples non-mandatory, and follows cross-clause dependencies before making a portability claim.\end{CornellReferenceOrientationBox}
\section{Self-Test Questions}
\begin{enumerate}
\item What role does Preamble play in the clause's overall model?
\item Which normative distinctions must be kept separate when applying Properties of expressions?
\item How would you audit a program or implementation that depends on Standard conversions?
\item Compare Usual arithmetic conversions with Primary expressions; what different question does each answer?
\item What role does Primary expressions play in the clause's overall model?
\item Which normative distinctions must be kept separate when applying Compound expressions?
\item How would you audit a program or implementation that depends on Constant expressions?
\item Compare Value category with Type; what different question does each answer?
\item What role does Type play in the clause's overall model?
\item Which normative distinctions must be kept separate when applying Context dependence?
\item How would you audit a program or implementation that depends on General?
\item Compare Lvalue-to-rvalue conversion with Array-to-pointer conversion; what different question does each answer?
\item What role does Array-to-pointer conversion play in the clause's overall model?
\item Which normative distinctions must be kept separate when applying Function-to-pointer conversion?
\item Scenario: a program relies on an unstated implementation behavior while using Temporary materialization conversion. What must be checked before calling the program portable?
\item Scenario: a program relies on an unstated implementation behavior while using Qualification conversions. What must be checked before calling the program portable?
\end{enumerate}
\clearpage\section{Answer Key}
\begin{enumerate}
\item Preamble is one part of the clause's core-language semantics model. Its role is understood by connecting the local rule in 7.1 to the clause orientation and its dependent concepts.
\item Keep formation and well-formedness separate from runtime preconditions and effects. Also distinguish mandatory wording from notes, implementation choice from undefined behavior, and interface declarations from detailed contracts; see 7.2.
\item Audit Standard conversions by locating the controlling wording in 7.3, listing inputs and type constraints, proving any preconditions, recording effects and failure behavior, and checking lifetime, invalidation, complexity, and synchronization where applicable.
\item Usual arithmetic conversions and Primary expressions occupy different steps or facility groups. Analyze each under its own subclause, then follow the explicit dependency between them rather than transferring one set of guarantees to the other; begin at 7.4.
\item Primary expressions is one part of the clause's core-language semantics model. Its role is understood by connecting the local rule in 7.5 to the clause orientation and its dependent concepts.
\item Keep formation and well-formedness separate from runtime preconditions and effects. Also distinguish mandatory wording from notes, implementation choice from undefined behavior, and interface declarations from detailed contracts; see 7.6.
\item Audit Constant expressions by locating the controlling wording in 7.7, listing inputs and type constraints, proving any preconditions, recording effects and failure behavior, and checking lifetime, invalidation, complexity, and synchronization where applicable.
\item Value category and Type occupy different steps or facility groups. Analyze each under its own subclause, then follow the explicit dependency between them rather than transferring one set of guarantees to the other; begin at 7.2.1.
\item Type is one part of the clause's core-language semantics model. Its role is understood by connecting the local rule in 7.2.2 to the clause orientation and its dependent concepts.
\item Keep formation and well-formedness separate from runtime preconditions and effects. Also distinguish mandatory wording from notes, implementation choice from undefined behavior, and interface declarations from detailed contracts; see 7.2.3.
\item Audit General by locating the controlling wording in 7.6.2.1, listing inputs and type constraints, proving any preconditions, recording effects and failure behavior, and checking lifetime, invalidation, complexity, and synchronization where applicable.
\item Lvalue-to-rvalue conversion and Array-to-pointer conversion occupy different steps or facility groups. Analyze each under its own subclause, then follow the explicit dependency between them rather than transferring one set of guarantees to the other; begin at 7.3.2.
\item Array-to-pointer conversion is one part of the clause's core-language semantics model. Its role is understood by connecting the local rule in 7.3.3 to the clause orientation and its dependent concepts.
\item Keep formation and well-formedness separate from runtime preconditions and effects. Also distinguish mandatory wording from notes, implementation choice from undefined behavior, and interface declarations from detailed contracts; see 7.3.4.
\item First identify the exact requirement in 7.3.5, then classify the relied-on behavior as required, unspecified, implementation-defined, conditionally supported, or undefined. Portability requires satisfying the program obligations and avoiding dependence on a choice not guaranteed in the target environments.
\item First identify the exact requirement in 7.3.6, then classify the relied-on behavior as required, unspecified, implementation-defined, conditionally supported, or undefined. Portability requires satisfying the program obligations and avoiding dependence on a choice not guaranteed in the target environments.
\end{enumerate}
\section{Key-Term Recap}
\begin{longtable}{@{}>{\RaggedRight\arraybackslash}p{0.28\textwidth}>{\RaggedRight\arraybackslash}p{0.64\textwidth}@{}}\toprule\textbf{Term} & \textbf{Working meaning}\\\midrule\endfirsthead\toprule\textbf{Term} & \textbf{Working meaning (continued)}\\\midrule\endhead\bottomrule\endfoot
value category & The glvalue/prvalue classification that governs identity, materialization, and binding.\\
standard conversion & An implicit conversion from one of the defined language conversion categories.\\
temporary materialization & Conversion of a prvalue into an xvalue denoting a temporary object.\\
sequenced before & An asymmetric execution relation between evaluations within a thread.\\
constant expression & An expression meeting the restrictions required for compile-time semantic use.\\
potentially evaluated & An expression context in which evaluation may contribute runtime effects.\\
discarded-value expression & An expression evaluated for effects while its value is discarded.\\
\end{longtable}
\CornellTakeaway{Expression correctness comes from the combined analysis of type, value category, conversions, lifetime, overload selection, and sequencing.}
\end{document}
